% Root exploratory note, written 19 November 2026 after the CP28 paper freeze.
% Outside the frozen paper. No main-gate movement is asserted.
\subsection*{Full hyperplane radicals and a dual alternating kernel map}

The statements below are elementary candidate lemmas with complete direct
proofs, awaiting independent audit. They supply conditional tests only.
No higher-degree Cartier tensor is identified with an alternating
three-form, and no previously withdrawn Pfaffian model is reused.

\paragraph{Full-radical generation.}
Let $V$ be an $n$-dimensional $\mathbf F_2$-space, $n\geq2$. For each
nonzero $\ell\in V^\vee$ put
\[
 U_\ell=\{1+v\ell:v\in\ker\ell\}\subseteq\operatorname{GL}(V).
\]
If $G$ is generated by a family of entire groups $U_\ell$ and $V^G=0$,
then $G=\operatorname{GL}(V)$.

\emph{Proof.} The intersection $K$ of the kernels of the selected
functionals is fixed pointwise by every generator, hence by $G$.
Thus $K=0$, so those functionals span $V^\vee$. Select a basis
$\ell_1,\ldots,\ell_n$ among them, with dual basis $e_1,\ldots,e_n$.
For $i\ne j$ the group $U_{\ell_i}$ contains $1+e_j\ell_i$.
These elementary row additions generate $\operatorname{GL}_n(\mathbf F_2)$:
Gaussian elimination uses row additions and row swaps, a swap being
three row additions in characteristic two, and there are no nontrivial
diagonal scalings. This proves the assertion.

The hypothesis is generation by the full radicals, not merely their
containment in $G$. If $N$ is their normal closure in a larger group,
one needs $V^N=0$, which does not follow from $V^G=0$. In the geometric
application, divisorial normal generation and fullness of the necessary
inertia groups both remain to be proved. This does not close the earlier
parabolic countermodel.

\paragraph{Differential of the left-kernel map.}
Let $k$ have characteristic two, let $V,W$ have the same finite dimension
$h$, and let $b:\Lambda^2V\to W$ be a linear map, regarded as an
alternating bilinear map. Set $M_s(t)=b(s,t)$. On the open subset
$\mathcal U\subset\mathbf P(V)$ where $M_s$ has rank $h-1$, its right
kernel is $ks$, and its left kernel determines a morphism
\[
 \kappa:\mathcal U\longrightarrow\mathbf P(W^\vee),
 \qquad [s]\longmapsto[\ell_s].
\]
For $\ell\in W^\vee$ define the alternating form
$B_\ell(v,w)=\ell(b(v,w))$. At every point of $\mathcal U$,
\[
             \operatorname{rank}(d\kappa_{[s]})
                      =\operatorname{rank}(B_{\ell_s}).
\]
Here differential rank is computed after a field extension defining the
point, and the equality is unchanged by rescaling $s$ or $\ell_s$.

\emph{Proof.} We have $\operatorname{im}M_s=\ker\ell_s$, so $M_s$
induces an isomorphism $V/ks\simeq\ker\ell_s$. Differentiate
$\ell_sM_s=0$ in the direction of $t$ to obtain
\[
             (d\ell_s(t))M_s+\ell_sM_t=0.
\]
Changing the representative of $d\ell_s(t)$ by a multiple of $\ell_s$
does not affect this equation. Restriction to $\ker\ell_s$ identifies
$W^\vee/k\ell_s$ with $(\ker\ell_s)^\vee$, and then $M_s$ identifies
this with $(V/ks)^\vee$. Under those identifications the derivative is
the map $t\mapsto B_{\ell_s}(t,-)$. Its source is $V/ks$ because
$s$ is in the radical of $B_{\ell_s}$. Quotienting by that radical
line does not change the rank of the form. This proves the equality.

In particular, if $h$ is odd and this rank is $h-1$ at one point,
$\kappa$ is separably dominant. In that case it is birational: on a
dense open its value $\ell$ has $\operatorname{rad}(B_\ell)$ equal
to a single line, which recovers $[s]$ rationally. The inverse can be
written by maximal Pfaffians of $B_\ell$; equivalently, ordinary
cofactor linear algebra gives the same unique radical line. No
Pfaffian formula for $M_s$ itself is required.

\paragraph{The rank inference that is not allowed.}
Generic rank $h-1$ of $M_s$ alone does not imply maximal differential
rank. For odd $h\geq5$ take $V=W$, choose $0\ne\ell_0\in V^\vee$,
and put
\[
                  b(s,t)=\ell_0(s)t+\ell_0(t)s.
\]
This is alternating. If $\ell_0(s)\ne0$, its right kernel is exactly
$ks$, its image is $\ker\ell_0$, and its rank is $h-1$.
Nevertheless $\kappa$ is the constant map $[\ell_0]$.
More generally $B_\ell=\ell_0\wedge\ell$ has rank at most two,
and $B_{\ell_0}=0$. This is an abstract bilinear countermodel, not
an asserted geometric Cartier realization.

\paragraph{Exact geometric status.}
The degree-one nonordinary principal-surface proof already establishes
a full radical of order $2^7$: its moving linearized kernel map on
the Hermitian cubic has differential rank seven. The root reread the
complete canonical uniform-cubic proof and the complete moving-kernel
Artin--Schreier lemma during this pass. For higher degree, the typed
identification of the relevant kernel map with the map above, its
maximal differential rank on an actual corank-one wall, and the
required global radical generation remain unproved. The possibility
of a Cartier bilinear realization is a research lead, not a theorem.
For a direct union of three odd-degree curves the marked-normalization
Lang subspace still has codimension four in the generic $f=0$ module
(two in the $f=1$ case); every presently constructed boundary operation
preserves it. Thus no odd-degree sum-closure theorem is inferred.

No computation was executed for this note. A literature pass searched
for alternating bilinear maps, left kernels, Pfaffian Cremona maps,
and Hasse--Witt/Cartier moving-kernel monodromy. The first engine gave
no directly applicable primary theorem; verification returned an
unopened secondary listing of \emph{On Birational Maps and Jacobian
Matrices}, Example 2.5, and unrelated material. No result from that
listing is used, no original proof from it was read, and no priority
is claimed. Exact searches are recorded in the companion log.
